% Part: first-order-logic
% Chapter: beyond
% Section: many-sorted-logic

\documentclass[../../../include/open-logic-section]{subfiles}

\begin{document}

\olfileid{fol}{byd}{msl}

\olsection{बहुवर्गीय तर्कशास्त्र}

प्रथम-क्रम तर्कशास्त्रात चरे आणि संख्यापक एकाच !!{domain} वर लागू होतात.
पण अनेक (परस्परविभक्त) !!{domain}s असणे बऱ्याचदा उपयुक्त ठरते: उदाहरणार्थ,
संख्यांचा !!a{domain}, भूमितीय वस्तूंचा !!a{domain}, संख्यांपासून संख्यांकडे
जाणाऱ्या फलनांचा !!a{domain}, आबेली गटांचा !!a{domain}, इत्यादी ठेवावेसे वाटू
शकतात.

बहुवर्गीय तर्कशास्त्र अशी चौकट पुरवते. सुरुवात ``वर्गांच्या'' यादीने होते---
वस्तूचा ``वर्ग'' ती कोणत्या ``!!{domain}'' मध्ये असणे अपेक्षित आहे ते दर्शवतो.
मग प्रत्येक वर्गासाठी स्वतंत्र !!{variable}s आणि संख्यापक, तसेच (सामान्यतः)
स्वतंत्र !!{identity} असते. फलने आणि संबंध कोणत्या वर्गांतील वस्तू फलसाधक
म्हणून घेऊ शकतात यानुसार तेही ``प्रकारबद्ध'' केलेले असतात. याखेरीज
प्रथम-क्रम तर्कशास्त्राचे नेहमीचे नियम तसेच ठेवले जातात; संख्यापकांचे नियम
प्रत्येक वर्गासाठी स्वतंत्रपणे पुनरावृत्त होतात.

उदाहरणार्थ, आंतरराष्ट्रीय संबंधांचा अभ्यास करण्यासाठी आपण फ्रेंच नागरिक आणि
जर्मन नागरिक असे वस्तूंचे दोन वर्ग असलेली भाषा निवडू शकतो. पहिल्या वर्गातील
वस्तूंसाठी ``द्राक्षारस पितो'' हा एकस्थानी संबंध; दुसऱ्या वर्गातील वस्तूंसाठी
``वुर्स्ट खातो'' हा दुसरा एकस्थानी संबंध; आणि दोन फलसाधक घेणारा ``बहुराष्ट्रीय
विवाहित जोडपे बनवतात'' हा द्विस्थानी संबंध असू शकतो, ज्याचा पहिला फलसाधक
पहिल्या वर्गातील आणि दुसरा फलसाधक दुसऱ्या वर्गातील असतो. फ्रेंच नागरिकांसाठी
$a$, $b$, $c$ ही चरे आणि जर्मन नागरिकांसाठी $x$, $y$, $z$ ही चरे वापरली, तर
\[
\lforall[a][\lforall x][(\Atom{\Obj{MarriedTo}}{a,x} \lif
(\Atom{\Obj{DrinksWine}}{a} \lor \lnot \Atom{\Obj{EatsWurst}}{x})]]
\]
हे सूत्र असे सांगते की कोणत्याही फ्रेंच व्यक्तीचा विवाह एखाद्या जर्मन
व्यक्तीशी झाला असेल, तर एकतर ती फ्रेंच व्यक्ती द्राक्षारस पिते किंवा ती जर्मन
व्यक्ती वुर्स्ट खात नाही.

बहुवर्गीय तर्कशास्त्र प्रथम-क्रम तर्कशास्त्रात स्वाभाविक रीतीने अंतःस्थापित
करता येते: बहुवर्गीय !!{domain}s मधील सर्व वस्तू एकत्र करून प्रथम-क्रमाचा एक
!!{domain} घ्यायचा, वर्गांची नोंद ठेवण्यासाठी एकस्थानी !!{predicate}s
वापरायची आणि संख्यापकांना योग्य वर्गापुरते सापेक्ष करायचे. उदाहरणार्थ, वरच्या
उदाहरणाशी संबंधित प्रथम-क्रम भाषेत वर्णिलेल्या इतर संबंधांबरोबर
``$\Obj{German}$'' आणि ``$\Obj{French}$'' ही एकस्थानी !!{predicate}s असतील;
मात्र वर्गाविषयीच्या अटी पुसलेल्या असतील. जर्मन वर्गातील !!a{variable}~$x$
असलेला वर्गबद्ध संख्यापक $\lforall[x][!A]$ पुढीलप्रमाणे रूपांतरित होतो:
\[
\lforall[x][(\Atom{\Obj{German}}{x} \lif !A)].
\]
वर्ग परस्परांपासून वेगळे आहेत याची खात्री देणारी स्वयंसिद्धकेही जोडावी
लागतात---उदाहरणार्थ,
$\lforall[x][\lnot (\Atom{\Obj{German}}{x} \land
  \Atom{\Obj{French}}{x})]$---तसेच ``द्राक्षारस पितो'' हा संबंध फक्त
$\Atom{\Obj{French}}{x}$ या विधेयाची पूर्ति करणाऱ्या वस्तूंनाच लागू होतो
याची हमी देणारी स्वयंसिद्धकेही जोडावी लागतात. या संकेतांमुळे आणि
स्वयंसिद्धकांमुळे बहुवर्गीय !!{sentence}s प्रथम-क्रम !!{sentence}s मध्ये आणि
बहुवर्गीय !!{derivation}s प्रथम-क्रम !!{derivation}s मध्ये रूपांतरित होतात,
हे दाखवणे अवघड नाही. तसेच बहुवर्गीय !!{structure}s संबंधित प्रथम-क्रम
!!{structure}s मध्ये आणि उलट दिशेनेही ``रूपांतरित'' होतात; म्हणून बहुवर्गीय
तर्कशास्त्रासाठीही आपल्याला संपूर्णता प्रमेय मिळते.

\end{document}
